% ===================== 练习 =====================
\phantomsection\label{sec:practice}
\hdr{\S11}{独立练习（五题）}

五题分别练：最低阶的定系数（P1）、Legendre 展开（P2）、球壳的两个径向条件（P3）、球对称的球贝塞尔（P4）、
带 $\varphi$ 依赖的球谐（P5）。题干只给条件与考点，步骤全在附录，题号对应。

\practice{P1}{自编，仿 §6.6 例 1}{半径 $a$ 的球内解 $\nabla^2u=0$，球面 $u(a,\theta)=U_0P_2(\cos\theta)$。求 $u(r,\theta)$，并指出球内 $|u|$ 的最大值出现在哪里。}{$l=2$ 的径向因子是 $r^2$；球内只留 $r^l$}{ans:P1}

\practice{P2}{自编，Legendre 展开}{把 $f(\theta)=\cos^3\theta$ 按 $P_l(\cos\theta)$ 展开，即求系数 $c_l$ 使 $f=\sum c_lP_l$。(1) 用 $P_1,P_3$ 的显式直接分解；(2) 用正交性公式验算 $c_1,c_3$，并说明为什么 $c_0=c_2=c_4=\cdots=0$。}{$x^3$ 用 $P_1,P_3$ 表示；奇偶性筛掉偶 $l$}{ans:P2}

\practice{P3}{自编，球壳}{球壳 $a<r<b$ 内解 $\nabla^2u=0$，内球面 $u(a,\theta)=0$，外球面 $u(b,\theta)=U_0\cos\theta$。求 $u(r,\theta)$，并核对两个球面条件。}{球壳两组径向解都留；$l=1$ 的两个系数由两个球面条件定}{ans:P3}

\practice{P4}{自编，球对称输运}{半径 $R$ 的球内，初始温度 $u(r,0)=U_0$（常数），球面保持零温 $u(R,t)=0$。求球对称解 $u(r,t)$ 的级数，并写出前两个非零项。}{$l=0$ 的球贝塞尔是 $j_0=\sin x/x$；$k_n=n\pi/R$；$\int_0^R j_0^2r^2dr$ 的模方}{ans:P4}

\practice{P5}{自编，带 $\varphi$ 依赖}{半径 $a$ 的球内解 $\nabla^2u=0$，球面 $u(a,\theta,\varphi)=U_0\sin\theta\cos\varphi$。求 $u$，并把它写成直角坐标形式。}{$P_1^1(\cos\theta)=\sin\theta$；$r\sin\theta\cos\varphi=x$}{ans:P5}

\clearpage
